\chapter{Elliptische Kurven}
Dieses Kapitel legt die Grundlagen für die Arbeit mit elliptischen Kurven. Neben der allgemeinen Gleichung wird auch die Gruppenstruktur eingeführt. Ein großer Punkt spielt dabei der Beweis der Assoziativitität.

\section{Grundlagen}
Das Kapitel basiert auf\cite{washington08} Washington 2008 S.9ff.

Beginnen wir also mit der Definition einer Elliptische Kurve.
\begin{definition}
    Eine elliptische Kurve über einem Körper K wird definiert mit
    \[
        y^2+a_1xy+a_3y=x^3+a_2x^2+a_4x+a_6,
    \]
    wobei wir das die Weierstraß-Gleichung nennen.
    Wenn nun gilt char($K\neq 2,3$), so können wir die kurze
    Weierstraß-Gleichung nutzen
    \[
        y^2=x^3+Ax+B.
    \]
    Wobei A und B konstante Werte sind.
\end{definition}
Wir schreiben für den Punkt im Unendlichen \(\infty =\mathcal{O}\). Die Punkte $L\subset K $ einer elliptische Kurve können dann mit
\[
    E(L)=\{\mathcal{O} \} \cup \{(x,y)\in L\times L\mid y^2=x^3+Ax+B\}
\]
dargestellt werden. Nach Definition ist der Punkt $\mathcal{O}$ immer enthalten, warum wir diesen benötigen wird im nächsten Kapitel genauer erklärt.
Nicht über jedem Körper lassen sich Elliptische Kurven sinnvoll optisch darstellen, daher betrachten wir zuerst den Körper der reellen Zahlen.
Hierbei gibt es zwei grundlegende Formen.

Über dem Körper \(\mathbb{R}\) kann die Kurve je nach Anzahl der reellen Nullstellen verschiedene Formen annehmemen. Zum Beispiel hat $y^2 = x^3 - 3x$ drei verschiedene reelle Nullstellen, nämlich $x = 0$ und $x = \pm \sqrt{3}$. Im Gegensatz dazu hat die Kurve $y^2 = x^3 + 3x$ lediglich eine reelle Nullstelle bei $x = 0$.\\
\begin{figure}[h]
    \centering
    \begin{subfigure}{0.35\linewidth}
        \centering
        \includegraphics[width=\linewidth]{kap2_bs1.png}
        \caption{Eigene Darstellung mit GeoGebra $y^2=x^3-3x$}
    \end{subfigure}
    \hfill
    \begin{subfigure}{0.45\linewidth}
        \centering
        \includegraphics[width=\linewidth]{kap2_bs2.png}
        \caption{Eigene Darstellung mit GeoGebra $y^2=x^3+3x$}
    \end{subfigure}
    \label{fig:vergleich}
\end{figure}

Wie wir in der Darstellung sehen können, ändert sich die Form der Kurve über dem Körper $\mathbb{R}$, wenn wir mehrere reelle Nullstellen haben. Das Polynom \(x^3+Ax+B\) darf daher keine mehrfachen Nullstellen besitzen, dies ist äquivalent dazu, dass die Kurve \emph{nicht singulär} ist.

\begin{definition}
    Sei $f \in K[x,y]$ und
    \[
        E = V(f) = \{(x,y)\in K^2 \mid f(x,y)=0\}
    \]
    eine affine ebene Kurve. Ein Punkt \((a,b)\in E\) heißt \textbf{singulär}, wenn,
    \[
        f_x(a,b)=0 \quad \text{und}\quad f_y(a,b)=0.
    \]
    Andernfalls heißt der Punkt \textbf{regulär}.

    Die Kurve \(E\) heißt \textbf{singulär}, wenn sie einen singulären Punkt besitzt.
    vgl. \cite{plaumann20} S.1,6
\end{definition}

Wir wollen nun, dass unsere elliptische Kurve auch diese Bedingung erfüllt.
Wir schreiben die Gleichung nun um zu
\[
    f(x,y)=y^2-x^3-Ax-B
\]
dann sind die Ableitung
\[
    f_x= -3x^2-A \quad \text{ und } \quad f_y=2y
\]
Es muss nun für beide ungleich 0 gelten,
\[
    f_x= -3x^2-A = 0 \Leftrightarrow 3x^2+A = 0
\]
und
\[
    f_y=2y = 0 \Leftrightarrow y = 0, \quad \text{char }\neq 2
\]
daraus folgt ein singulärer Punkt muss die Form $(x,0)$ haben und zusätzlich muss \(x^3+Ax+B = 0\) gelten, sodass der Punkt auf der Kurve liegt.
Also x darf nicht eine Nullstelle von \(x^3+Ax+B\) und der Ableitung \(3x^2+A\) sein.
Die Diskriminante für ein kubisches Polynom
\[
    p(x)=x^3+Ax+B
\]
ist hierbei
\[
    \Delta  = -4A^3-27B^2
\]
(\cite{karpf24} vgl. Karpf 2024, S.470f.).
Damit besitzt $p$ genau dann keine mehfachen Nullstellen, wenn
\[
    4A^3+27B^2 \neq 0.
\]
Für die Nullstellen $r_1, r_2, r_3$ des Polynoms gilt
\[
    ((r_1 - r_2)(r_1 - r_3)(r_2 - r_3))^2 = -(4A^3 + 27B^2).
\]
Insbesondere folgt daraus, dass die Nullstellen genau dann paarweise verschieden sind, wenn die Diskriminante $\Delta = -(4A^3 + 27B^2)$ ungleich Null ist.
Im vorliegenden Fall ist $\Delta < 0$, sodass voneinader verschiedene reelle Nullstellen vorliegen.


\section{Gruppenstruktur}

Die Darstellung der Gruppenstruktur elliptischer Kurven orientiert sich in diesem Abschnitt an \cite{washington08} vgl. Washington 2008 S.12ff. und \cite{silverman06} Silverman 2006 S.8-19. Wobei der Beweis der Wohldefiniertheit der Addition auf \cite{zwegers24} vgl. Zwegers 2024 S. 477-481 basiert, dieser wird jedoch auf die bereits kurze Weierstraßform abgewandelt.

Wir beginnen mit der Addition von zwei Punkten auf einer elliptischen Kurve. Dazu seien $P=(x_1,y_1)$ und $Q=(x_2,y_2)$ zwei Punkte auf einer elliptischen Kurve $E\colon y^2=x^3+Ax+B$. Gesucht ist der Punkte $R= P+_EQ$. Für die graphische Darstellung wählen wir ein Beispiel über $\mathbb{R}$ mit der Kurve $y^2=x^3-4x+10$.

Im ersten Schritt wird die Gerade $L$ durch $P$ und $Q$ gezeichnet.
\begin{figure}[htbp]
    \centering
    \includegraphics[width=0.5\linewidth]{kap2_2_bsp1.png}
    \caption{Gerade $L$ durch  $P$ und $Q$}
    \label{fig:punktaddition-gerade}
\end{figure}

Als zweiten Schritt nehmen wir den Schnittpunkt $R'$ von der Gerade $L$ und $E$. Anschließend spiegeln wir $R'$ an der x-Achse und erhalten dadurch $R$. Wir definieren $R=P+_E Q$. Hierbei ist wichtig hervorzuheben, dass es sich dabei von der Addition der Koordinaten zweier Punkte unterscheidet.

\begin{figure}[htbp]
    \centering
    \includegraphics[width=0.5\linewidth]{kap2_2_bsp2.png}
    \caption{Schnittpunkt $R'$ von Gerade $L$ und Kurve $E$ und die Spiegelung $R$}
    \label{fig:punktaddition-dritter-sp}
\end{figure}

Nun betrachten müssen wir noch den Fall betrachten, wenn wir einen Punkt $P$ auf sich selbst addieren wollen auf der Kurve $E$. Da es keinen zweiten Punkt, sodass durch beide Punkte eine Gerade gezogen werden kann, nehmen wir an der zweite Punkt $Q$ nähert sich P an, sodass die Gerade $L$ zu einer Tangente von $E$ bei $P$ wird.
\begin{figure}[htbp]
    \centering
    \includegraphics[width=0.5\linewidth]{kap2_2_bsp3.png}
    \caption{Tangente am Punkt $P$}
    \label{fig:punktaddition-spiegelung}
\end{figure}

Anschließend nehmen wir wieder den Schnittpunkt $R'$ von der Kurve $E$ und der Tangente $L$ und spiegeln diesen an der x-Achse. Dadurch erhalten wir wieder den Punkt $P+_EP=R$, denn wir aber nun als $2P$ bezeichnen.

\begin{figure}[htbp]
    \centering
    \includegraphics[width=0.6\linewidth]{kap2_2_bsp4.png}
    \caption{Spiegelung von $R'$, wodurch $R=2P$ entsteht}
    \label{fig:punktverdoppelung-tangente}
\end{figure}

Als letzten Fall schauen wir uns noch an, wenn die Gerade $L$ keinen Schnittpunkt mit $E$ hat. Hierfür betrachten wir wieder den Punkt $P$ und seinen an der x-Achse gespiegelten Punkt $P'=-P$.
\begin{figure}[htbp]
    \centering
    \includegraphics[width=0.5\linewidth]{kap2_2_bsp5.png}
    \caption{Schnittpunk von L im Punkt $\mathcal{O}$}
    \label{fig:punktverdopplung-ergebnis}
\end{figure}
Dafür brauchen wir einen dritten Punkte um $P+_E(-P)$ definieren zu kennen wir kreieren dafür einen Punkt $\mathcal{O}$. Dieser Punkt ist dabei auf jeder vertikalen Gerade

\subsection*{Herleitung der Additionsformeln}
Die folgende Herleitung orientiert sich an\cite{washington08} Washington 2008, S. 13-15,\cite{kapl05} vgl. Kaplan 2005 S.179ff. und\cite{buell24} Buell 2024, S.99.

Seien nun $P=(x_1,y_1)$ und $Q=(x_2,y_2)$ Punkte auf $E$ mit $P,Q\neq \mathcal{O}$.
Zuerst betrachten wir den Fall $P\neq Q$ und \(x_1\neq x_2\). Dann hat die Gerade \(L\) durch durch \(P\) und \(Q\) die Form
\[
    y=m(x-x_1)+y_1
    \qquad \text{mit} \qquad
    m=\frac{y_2-y_1}{x_2-x_1}.
\]
Zur Bestimmung des dritten Schnittpunkts setzen wir die Geradengleichung in die Kurvengleichung
\[
    E: y^2=x^3+Ax+B
\]
ein. Es gilt also
\[
    \bigl(m(x-x_1)+y_1\bigr)^2=x^3+Ax+B.
\]
Nach Ausmultiplizieren und Umstellen erhält man
\[
    0=x^3-m^2x^2+\bigl(2m^2x_1-2my_1+A\bigr)x
    +\bigl(B-m^2x_1^2+2mx_1y_1-y_1^2\bigr).
\]

Die Nullstellen des kubischen Polynoms sind die \(x\)-Koordinaten der drei Schnittpunkte von \(L\) mit \(E\), also \(x_1\), \(x_2\) und \(x_3'\). Daher gilt
\[
    0=(x-x_1)(x-x_2)(x-x_3')
    =x^3-(x_1+x_2+x_3')x^2+(x_1x_2+x_1x_3+x_2x_3)x-x_1x_2x_3
\]
Mit Kooeffizientenvergleich erhält man
\[
    x_1+x_2+x_3' = m^2 \Leftrightarrow  x_3' = m^2-x_1-x_2
\]
Da \(R'=(x_3',y_3')\) auf der Geraden \(L\) liegt, gilt
\[
    y_3'=m(x_3'-x_1)+y_1.
\]
Die Summe \(P+Q\) ist die Spiegelung von \(R'\) an der \(x\)-Achse. Somit ist
\[
    P+Q=(x_3,y_3)
\]
mit
\[
    x_3=x_3'=m^2-x_1-x_2,
    \qquad
    y_3=-y_3'=m(x_1-x_3)-y_1.
\]

\textbf{voll restliche Herleitung Einfügen, Rosen S.12-14}
\begin{definition}[Punktaddition]
    \label{def:punktaddition}
    Die Punktaddition für einen Punkt $P+Q=R=(x_3,y_3)$ definieren wir wie folgt:
    \begin{enumerate}
        \item Wenn $x_1\neq x_2$, dann
              \[x_3=m^2-x_1-x_2, \quad y_3=m(x_1-x_3)-y_1 \quad \text{ mit  } m=\frac{y_2-y_1}{x_2-x_1}.\]
        \item Wenn $x_1=x_2$, aber $y_1\neq y_2$, dann ist
              \[P+Q=\mathcal{O}.\]
        \item Wenn $P=Q$ und $y_1\neq 0$, dann
              \[x_3=m^2-2x_1, \quad y_3=m(x_1-x_3)-y_1, \quad \text{ mit } m=\frac{3x_1^2+A}{2y_1}.\]
        \item Wenn $P=Q$ und $y_1=0$ dann,
              \[P+Q=\mathcal{O}.\]
    \end{enumerate}

    Zusätzlich definieren wir
    \[P+\mathcal{O}=P\]
    für alle Punkte $P$ auf $E$.
\end{definition}


\begin{lemma}[Wohldefiniertheit der Punktadditon]
    \label{lem:wohldefiniertheit}
    Sei \(P_1=(x_1,y_1)\) und \(P_2=(x_2,y_2)\) in $E\setminus \{\mathcal{O}\}$ mit \((x_2,y_2)\neq (x_1,-y_1)\). Dann gilt
    \[
        x_3 = m_{12}^2-x_1-x_2 \quad \text{ und } \quad y_3=-y_1-m_{12}(x_3-x_1)
    \]
    mit
    \[
        m_{12}=m(P_1,P_2)=
        \begin{cases}
            \dfrac{y_2-y_1}{x_2-x_1}, & \text{wenn } x_1 \neq x_2, \\[6pt]
            \dfrac{3x_1^2+A}{2y_1},   & \text{wenn } P_1=P_2.
        \end{cases}
    \]
    Dann gilt $(x_3,y_3)\in E\setminus \{\mathcal{O}\}$
\end{lemma}

Damit ist die oben definierte Punktaddition abgeschlossen auf \(E\).

\begin{proof}
    Also es gilt zu zeigen, dass gilt $P_1,P_2\in E \Rightarrow P_1 +_E P_2 \in E$. Also wenn $x_1=x_2$, dann gilt \(f(x_1)=y_1^2 \) und \(f(x_2)=y_2^2\), also $y_1^2=y_2^2$, damit gilt \(y_1= y_2 \) oder \(y_1=-y_2\). Da wir bereits definiert haben, dass $(x_2,y_2)\neq (x_1,-y_1)$ gilt im Lemma, folgt daraus $y_1=y_2$ und damit $P_1=P_2$. Wäre $y_1=0$, so wäre wegen $y_1=y_2$ auch $y_2=0=-y_1$. Damit wäre es im Widerspruch zur Voraussetzung des Lemmas. Somit gilt $y_1\neq 0$ und damit $m_{12}\in K$ und $(x_3,y_3)\in K^2$ für diesen Fall wohldefiniert.
    Zur einfacher Notation $f(x)=x^3+Ax+B$, also die rechte Seite der Kurvengleichung. Wir führen nun das Polynom ein
    \[
        g(x)=f(x)-(y_1+m_{12}(x-x_1))^2\in K[x].
    \]
    Da $(x_1,y_1)$ auf unserer Kurve liegt, folgt daraus $g(x_1)=f(x_1)-(y_1+m_{12}(x_1-x_1))^2=f(x_1)-y_1^2=0$. Damit folgt daraus $x_1$ ist eine Nullstelle von $g(x)$.
    Sei nun $x_2\neq x_1$, dann folgt $y_1+m_{12}(x_2-x_1)=y_2$. Damit gilt dann $g(x_2)=f(x_2)-(y_1+\frac{y_2-y_1}{x_2-x_1}(x_2-x_1))^2 =f(x_2)-y_2^2$ und damit ist $x_2$ auch eine Nullstelle.

    Nun betrachten wir noch den Fall $x_1=x_2$, die Ableitung von $g$ ist dann,
    \[
        g'(x)=f'(x)-2m_{12}(y_1+m_{12}(x-x_1)) \in K[x].
    \]
    Für $x_1$ gilt dann, \(g'(x_1)=f'(x_1)-2m_{12}(y_1)= (3x_1^2+A)-2\frac{3x_1^2+A}{2y_1}y_1=0 \), damit ist dann $x_1$ eine Nullste von $g'(x)$ und damit dann sogar eine doppelte Nullstelle von $g(x)$.
    Wir können also in beiden Fällen schreiben
    \[
        g(x)=(x-x_1)(x-x_2)h(x), \quad h(x)\in K[x]
    \]
    Da $g(x)$ ein normiertes Poylnom vom Grad 3 ist, können wir $h(x)$ umschreiben zu $h(x)=x-x_0$ mit $x_0\in K$. Daraus folgt dann
    \begin{equation}
        \label{eq:faktorisierung-wohldefiniertheit}
        f(x)-\left(y_1+m_{12}(x-x_1)\right)^2
        =(x-x_1)(x-x_2)(x-x_3).
    \end{equation}
    Durch Vergleich der Koeffizienten von $x^2$ erhalten wir auf der linken Seite den Koeffizienten
    \[
        -m_{12}^2
    \]
    und auf der rechten Seite den Koeffizienten
    \[
        -(x_1+x_2+x_0).
    \]
    Damit gilt
    \[
        -m_{12}^2=-(x_1+x_2+x_0).
    \]
    Daraus folgt
    \[
        m_{12}^2=x_1+x_2+x_0
    \]
    und somit
    \[
        x_0=m_{12}^2-x_1-x_2=x_3.
    \]
    und daraus folgt dann $x_0=x_3$. Da $x_0$ eine Nullstelle von $g$ ist, gilt
    \[
        g(x_0)
        =
        f(x_0)
        -
        \left(y_1+m_{12}(x_0-x_1)\right)^2
        =0.
    \]
    Folglich gilt
    \[
        f(x_0)
        =
        \left(y_1+m_{12}(x_0-x_1)\right)^2.
    \]

    Nach Definition der Additionsformel ist
    \[
        y_3
        =
        -y_1-m_{12}(x_3-x_1)
        =
        -\left(y_1+m_{12}(x_3-x_1)\right).
    \]
    Durch Quadrieren erhalten wir
    \[
        y_3^2
        =
        \left(y_1+m_{12}(x_3-x_1)\right)^2.
    \]
    Da $x_0=x_3$ gilt, folgt
    \[
        y_3^2
        =
        \left(y_1+m_{12}(x_0-x_1)\right)^2
        =
        f(x_0)
        =
        f(x_3).
    \]
    Damit erfüllt $(x_3,y_3)$ die Kurvengleichung und somit gilt
    \[
        (x_3,y_3)\in E(K).
    \]
\end{proof}
\begin{bemerkung}
    \label{bem:sum-ist-identitaet}
    Damit ist der nichttriviale Fall der Punktaddition wohldefiniert.
    Für $P_2=-P_1$ gilt nach Definition
    \[
        P_1+P_2=\mathcal{O}\in E,
    \]
    und für $P\in E$ gilt
    \[
        P+\mathcal{O}=\mathcal{O}+P=P\in E.
    \]
    Somit ist die Punktaddition insgesamt abgeschlossen auf $E$.
\end{bemerkung}

\begin{theorem}
    \label{theo:eigenschaften-addition}
    Die Menge $E$ bildet bezüglich der Punktaddition eine abelsche Gruppe
    mit neutralem Element $\mathcal{O}$. Insbesondere gelten für alle
    $P,Q,R\in E$:
    \begin{enumerate}
        \item [Existenz der Identität:] \(P+\mathcal{O}=\mathcal{O}+P=P\) für alle \(P\in E\)
        \item [Kommutativität:] \(P+Q=Q+P\) für alle \(P,Q\in E\)
        \item [Existenz des Inversen:] Sei $P\in E$ gegeben, dann exestiert ein $-P\in E$ mit \(P+(-P)= \mathcal{O}\), für alle $P\in E$.
        \item [Assoziativität:] \((P+Q)+R=P+(Q+R)\) für alle \(P,Q,R\in E\)
    \end{enumerate}
\end{theorem}

\begin{proof}
    Die Existenz der Identität ist durch die Definition der Elliptische Kurve gegeben. Die Kommutativität kann einfach nachgerechnet werden indem die oben definieren Regeln für die Addition zweier Punkte anwenden.
    Für  das Inverse müssen wir von einem Punkt $P$ sein Reflektion an der x-Achse $P'$ nehmen. Dann erhalten wir $P+P'=\mathcal{O}.$
    Der Beweis der Assoziativität ist recht langwierig, weshalb er in \ref{sec:assoziativitaet} ausführlich behandelt wird.
\end{proof}


\begin{beispiel}
    Wir betrachten dafür eine Elliptische Kurve $E$ mit
    \[y^2=x^3-4x+4,\]
    auf dieser Kurve liegen zum Beispiel die Punkte $P=(0,2)$ und $Q=(2,2)$. Nun berechnen wir auf den vorherigen Regeln, da $x_1\neq x_2$ folgt,
    \[
        m = \frac{2-2}{2-0}=0
    \]
    \[
        x_3 = (0)^2-0-2=-2, \quad  y_3 = 0\cdot(0-(-2))-2 =-2
    \]
    also
    \[
        P+_EQ=(0,0)+_E(0,2)= (-2,-2)
    \]
    Analog können wir beliebige Punkte auf der Kurve addieren.
    \[(0,2)+_E(-2,-2)=(6,14)\]
\end{beispiel}

\begin{beispiel}
    Nun können wir auch Punkte berechnen, wenn wir diese geometrisch nicht sinnvoll darstellen können. Die Kurve \(y^2=x^3-4x+4\pmod 7\) also über dem endlichen Körper $\mathbb{F}_7$. Als Errinnerung die Kurve muss über dem Körper Singulär sein, es muss also \(4A^3+27B^2=4\cdot(-4)^3+27\cdot4^2=176\equiv1\pmod 7\neq0\) gelten.

    Wir wählen nun die Punkte $P=(0,2)$ und \(Q=(5,5)\)
    auf $E(\mathbb{F}_7)$.
    Es gilt nun
    \[
        m=\frac{5-2}{5-0}=\frac{3}{5} = 3\cdot 5^{-1}\equiv 3\cdot 3 \pmod 7= 9 = 2 \pmod 7
    \]
    und damit
    \[
        x_3=2^2-0-5=-1\equiv 6 \pmod 7, \quad y_3=2(0-6)-2=-14 \equiv 0\pmod 7
    \]
    also
    \[P+_EQ=(6,0).\]
\end{beispiel}

\subsection*{Skalar mal Punkt}(The aglebra of Elliptic Curves, S.30f.)
Um später Diffie-Hellman über der endlcihen Gruppe $E(\mathbb{F}_p)$ anwenden zu können(wird im Kryptographie Kapitel genauer erklärt), muss man auch für größere $p$ berechnen können, also
\[mP=P+P+\dots+P\in E(\mathbb{F}_p)\] für sehr große Werte von m.
Lass \(m\in\mathbb{N}_{>0}\).

Wir schreiben zuerst
\[m=m_0+m_1\cdot2+m_2\cdot 2^2+\dots+m_r\cdot 2^r, \quad m_0,\dots,m_r\in\{0,1\}\]
also zerlegen m in seinen einzelnen Zweierpotenzen. Dann kann $mP$ berechnet werden als
\[mP=m_0P+m_1\cdot 2P+m_2\cdot2^2P+\dots +m_r\cdot 2^rP,\]
wobei \(2^kP=2\cdot 2\cdots 2P\) benötigt nur k Doppelungen und damit können wir $mP$ in $O(\text{log } m)$ berechnen. Die Binärdarstellung von $m$ hat
\[
    r+1 = \lfloor\log_2m \rfloor+1
\].
Im Schnitt braucht man also $\log_2(m)$ Verdoppelungen und \(\frac{1}{2}\log_2(m)\) Additionen um $mP$ zu berechnen.

Man kann die benötigte Zeit nochmal reduzieren, indem man die Non-Adjacent Form verwendet. Die Addition und Subtraktion eines Punktes kosten genau gleich viel Zeit, also schreiben wir
\[
    m=m_0+m_1\cdot 2+m_2\cdot 2^2+\cdots+m_r\cdot 2^r, \quad m_0,\dots,m_r\in\{-1,0,1\}.
\]
Hierfür können wir im Schnitt  $\frac{2}{3}$ der $m_i=0$ für die Zerlegung nutzen, dass reduziert die Anzahl der durchschnittlichen Additionen auf \(frac{1}{3}\log_2(m)
\).
Aufgrund dieses Vorgehens ist die Berechnung von $m$ aus $P$ und $mP$ sehr schwierig. Das sogennante \ref{sec:diffie-hellman} für Elliptische Kurven ist später essenziell für das kryptographische Verfahren, das auf diesen aufbaut.


\section{Projektiver Raum und der Punkt bei Unendlich} (vgl. \cite{washington08}, S. 18ff.)
(vgl. \cite{clader25}, S.245f.)

Die Gruppenstruktur elliptischer Kurven wurde zuvor über die geometrische
Punktaddition beschrieben. Dabei spielt der Punkt im Unendlichen eine besondere
Rolle, da er als neutrales Element der Gruppe dient. Um diesen Punkt geometrisch
präzise zu beschreiben, gehen wir nun von der affinen Ebene zur projektiven Ebene über.

\begin{definition}[Projektive Ebene]
    Sei \(K\) ein Körper. Die projektive Ebene über \(K\) ist definiert als der Quotient
    \[
        \mathbb P^2_K
        =
        \left(K^3\setminus0\right)/\sim,
    \]
    wobei zwei Tripel \((x,y,z),(x',y',z')\in K^3\setminus\{(0,0,0)\}\) genau dann in einer Relation  stehen \(
    (x,y,z)\sim (x',y',z)\), wenn ein \(\lambda\in K^\times:
    (x,y,z)=\lambda(x',y',z')\) exestiert.
    Die Äquivalenzklasse von \((x,y,z)\) wird mit
    \((x:y:z)\) bezeichnet und heißt Punkt der projektiven Ebene.
\end{definition}
Die Koordinaten \((x:y:z)\) werden homogene Koordinaten genannt. Dabei bezeichnet \((x:y:z)\) nicht einen einzelnen Punkte, sondern die Äquivalenzklasse aller von Null verschiedenen skalaren Vielfachen von \((x,y,z)\). Für jedes \(\lambda\in K^\times\) gilt also
\[
    (x:y:z)=(\lambda x:\lambda y:\lambda z).
\]
Die Menge
\[
    g_\mathcal{O}=\{(x:y:z)\in \mathbb P^2_K\mid z=0\}.
\]
heißt die unendlich ferne Gerade der projektiven Ebene (vgl. \cite{forster15}, S.200ff.).
Ist $z \neq 0$, so kann z auf 1 normiert werden, denn es gilt \((x:y:z)=(x/z:y/z:1)\). Nun kann man eine Abbildung definieren
\[
    K^2\longrightarrow\mathbb{P}^2_K, \quad
    (x,y)\mapsto (x:y:1).
\]
Diese Punkte entsprechen genau den affinen Punkten \((x,y)\in K^2\) die bijektiv auf \(\mathbb{P}^2_K\setminus g_\mathcal{O}\) abgebildet werden.


\begin{lemma}
    Die Punkte auf der Geraden im Unendlichen \(L_\mathcal{O}\) entsprechen den Richtungen affiner Geraden. Insbesondere schneiden sich parallel affine Geraden in der projektiven Ebene in demselben Punkt auf \(L_\mathcal{O}\).
\end{lemma}

\section{Beweis der Assoziativitität}
\label{sec:assoziativitaet}
Wie bereits angekündigt benötigt der Beweis der Assoziativität ein eignes Unterkapitel, der Beweis basiert dabei neben der Quelle von Washington auf dem Paper von \cite{zwegers24} vgl. Zweger S. 477-486.


\begin{theorem}
    Seien \(P_1,P_2,P_3\in E\), dann gilt \((P_1+P_2)+P_3=P_1+(P_2+P_3)\).
\end{theorem}

Doch bevor wir die Assoziativitität beweisen können benötigen wir eine Reihe von Hilfsaussagen.

\begin{lemma}
    \label{lem:dreipunktsumme}
    Sei $P_1,P_2,P_3\in E\setminus{\{\mathcal{O}\}}$ mit $P_3\neq \pm P_1, P_2\neq -P_1, P_2\neq -P_3$ und $P_1+P_2\neq -P_3$. Dann werden die Koordinanten von Punkte $P_4=(P_1+P_2)+P_3$ beschrieben durch

    Dann besitzt der Punkt \(P_4=(P_1+P_2)+P_3\) die Koordinaten
    \[
        x_4=-x_1+x_2-x_3+\frac{1}{c_2^2}-2\frac{c_1}{c_2}, \quad y_4 = -y_2-(x_4-x_2)(c_1+x_4c_2)
    \]
    wobei
    \[
        c_1 := \frac{x_1m(P_2,P_3)-x_3m(P_1P_2)}{x_1-x_3}, \quad c_2:=\frac{m(P_1,P_2)-m(P_2,P_3)}{x_1-x_3}.
    \]
\end{lemma}

\begin{hilfssatz}
    \label{hilfs:fixpunkt-addition}
    Sei $P_1,P_2\in E$ mit $P_1+P_2=P_1$. Dann gilt $P_2 =\mathcal{O}$.
\end{hilfssatz}

\begin{proof}
    Wenn $P_1=\mathcal{O}$, dann gilt \(P_2=\mathcal{O}+P_2=P_1+P_2=P_1=\mathcal{O}\) und damit ist die Bedingung erfüllt.
    Wenn nun $P_2=-P_1$ gilt, dann gilt für die Bedingung nach Bemerkung \ref{bem:sum-ist-identitaet} $P_1 = P_1+P_2 = \mathcal{O}$ und damit $P_2=-P_1=-\mathcal{O}=\mathcal{O}$. Es gilt nun zu zeigen, dass es kein $P_1,P_2\in E\setminus\{\mathcal{O}\}$ mit $P_2\neq -P_1$ und $P_1+P_2=P_1$ exestiert.
    Angenommen es gibt einen solchen Fall dann erhalten wir mit der Wohldefiniertheit der Additon~\ref{lem:wohldefiniertheit}: $(x_1,y_1)=(x_3,y_3)$ und damit auch $y_1=y_3=-y_1-m_{12}(x_3-x_1)=-y_1-m_{12}\cdot 0 = -y_1 \Leftrightarrow y_1=0$. Wir nehmen nun wieder die Gleichung~\eqref{eq:faktorisierung-wohldefiniertheit}
    \[
        f(x)-\left(y_1+m_{12}(x-x_1)\right)^2=
        f(x)-m_{12}^2(x-x_1)^2 = (x-x_1)(x-x_2)(x-x_3) = (x-x_1)^2(x-x_2).
    \]
    Damit wäre aber $x_1$ eine doppelte Nullstelle von unserer elliptischen Kurve $f(x)$ und das Widerspricht der Annahme, dass sie singulär sein muss. Daraus folgt eine solche Annahme kann nicht gelten.
\end{proof}

\begin{hilfssatz}
    \label{hilfs:steigungsidentitaet}
    Sei $P_1,P_2,P_3\in E\setminus\{\mathcal{O}\}$ mit $P_3 \neq \pm P_1,P_2\neq -P_1,P_2\neq -P_3$ und $P_1+P_2\neq -P_3$. Außerdem ist $c_2$ definiert wie in Lemma~\ref{lem:dreipunktsumme}. Dann gilt
    \[
        c_2(m_{12}+m_{12,3})=1.
    \]

\end{hilfssatz}

\begin{proof}
    Es gilt wie beim Beweis der Wohldefiniertheit der Addition~\ref{lem:wohldefiniertheit}
    ist nach Gleichung~\eqref{eq:faktorisierung-wohldefiniertheit} bekannt, dass
    \[
        f(x)-\left(y_1+m_{12}(x-x_1)\right)^2
        =(x-x_1)(x-x_2)(x-\hat{x}).
    \]
    Wenn $P_2\neq P_3$, dann gilt
    \[
        c_2=\frac{m(P_1,P_2)-m(P_2,P_3)}{x_1-x_3}=\frac{m_{12}-\frac{y_3-y_2}{x_3-x_2}}{x_1-x_3}=\frac{\frac{m_{12}(x_3-x_2)}{x_3-x_2}-\frac{y_3-y_2}{x_3-x_2}}{x_1-x_3}
        = \frac{-y_3+y_2+m_{12}(x_3-x_2)}{(x_1-x_3)(x_3-x_2)}.
    \]
    Wir nutzen dafür, dass $y_2-y_1=m_{12}(x_2-x_1)$ gilt für die Gerade durch die Punkte $P_1$ und $P_2$. Dies formen wir um zu $y_2=y_1+m_{12}(x_2-x_1)$ und setzen es in die Gleichung ein
    \[
        \frac{-y_3+(y_1+m_{12}(x_2-x_1))+m_{12}(x_3-x_2)}{(x_1-x_3)(x_3-x_2)}
        = \frac{-y_3+y_1+m_{12}(x_3-x_1)}{(x_1-x_3)(x_3-x_2)}
    \]
    Wenn nun $P_1+P_2\neq P_3$ gilt. Wir definieren $(x_1,y_1) + (x_2,y_2) = (\hat{x},\hat{y})$ und rechnen
    \[
        m_{12}+m_{12,3}=m_{12}+\frac{y_3-\hat{x}}{x_3-\hat{x}}=m_{12}+\frac{y_3+y_1+m_{12}(\hat{x}-x_1)}{x_3-\hat{x}}
    \]
    \[
        =\frac{m_{12}(x_3-\hat{x})}{x_3-\hat{x}}+\frac{y_3+y_1+m_{12}(\hat{x}-x_1)}{x_3-\hat{x}}
    \]
    \[
        =\frac{y_3+y_1+m_{12}(x_3-x_1)}{x_3-\hat{x}}
    \]
    Wir betrachten nun 4 Fälle.
    \begin{enumerate}
        \item Für $P_2\neq P_3$ und $P_1+P_2\neq P_3$
              \[
                  c_2(m_{12}+m_{12,3}) = \frac{-y_3+y_1+m_{12}(x_3-x_1)}{(x_1-x_3)(x_3-x_2)} \cdot \frac{y_3+y_1+m_{12}(x_3-x_1)}{x_3-\hat{x}}
              \]
              \[
                  = \frac{(-y_3+y_1+m_{12}(x_3-x_1))(y_3+y_1+m_{12}(x_3-x_1))}{(x_1-x_3)(x_3-x_2)(x_3-\hat{x})}
              \]
              \[
                  = \frac{-y_3^2+m_{12}^2(x_3-x_1)^2+2\cdot m_{12}(x_3-x_1)y_1+ y_1^2}{(x_1-x_3)(x_3-x_2)(x_3-\hat{x})}
              \]
              \[
                  = \frac{-y_3^2+(y_1+ m_{12}(x_3-x_1))^2}{(x_1-x_3)(x_3-x_2)(x_3-\hat{x})}
              \]
              Es gilt $y_3^2=x^3+Ax+B$ mit $x=x_3$, dann $y_3^2=x_3^3+Ax_3+B$. Mit ~\eqref{eq:faktorisierung-wohldefiniertheit} gilt dann
              \[
                  \frac{-(f(x_3)-(y_1+m_{12}(x_3-x_1))^2)}{(x_1-x_3)(x_3-x_2)(x_3-\hat{x})}
              \]
              \[
                  \frac{-(x_3-x_1)(x_3-x_2)(x_3-\hat{x})}{-(x_3-x_)(x_3-x_2)(x_3-\hat{x})}= 1
              \]
        \item Sei nun $P_2\neq P_3$ und $P_1+P_2=P_3$ es gilt erneut $(\hat{x},\hat{y})=(x_3,y_3)$ mit $y_3=-y_1-m_{12}(x_3-x_1)$
              \[
                  c_2 = \frac{-y_3+y_2+m_{12}(x_3-x_2)}{(x_1-x_3)(x_3-x_2)} = \frac{-2y_3}{(x_1-x_3)(x_3-x_2)}
              \]
              Wir leiten nun die Annahme von ~\eqref{eq:faktorisierung-wohldefiniertheit} ab und setzen dann $x=x_3=\hat{x}$
              \[
                  f'(x)-2m_{12}(y_1+m_{12}(x-x_1))=f'(x_3)-2m_{12}(y_1+m_{12}(x_3-x_1))
              \]
              Es gilt dann $\hat{y}=-y_1-m_{12}(\hat{x}-x_1)$ und $\hat{y}=y_3$ und damit $y_1+m_{12}(x_3-x_1)=-y_3$
              \[
                  f'(x_3)+2m_{12}y_3
              \]
              Recht Seite abgeleitet ergibt
              \[
                  (x-x_2)(x-\hat{x})+(x-x_1)(x-\hat{x})+(x-x_1)(x-x_2)
              \]
              wieder $x=\hat{x}=x_3$ einsetzen
              \[
                  (x_3-x_2)(\hat{x}-\hat{x})+(x_3-x_1)(\hat{x}-\hat{x})+(x_3-x_1)(\hat{x}-x_2)
                  = (x_3-x_1)(\hat{x}-x_2)
              \]
              Also Gesamt ergibt sich
              \[
                  f'(x_3)+2m_{12}y_3 = (x_3-x_1)(\hat{x}-x_2)
              \]
              umformen für $y_3\neq 0$ gilt aus der Bedingung von Wohldefiniertheit
              \[
                  \frac{f'(x_3)}{2y_3}+m_{12} = \frac{(x_3-x_1)(\hat{x}-x_2)}{2y_3}
              \]
              Es ist $m_{12,3}=\frac{f'(x_3)}{2y_3}$ und daraus folgt
              \[
                  m_{12}+m_{12,3}= \frac{(x_3-x_1)(\hat{x}-x_2)}{2y_3}
              \]
              Damit gilt dann offensichtlich mit der Vertauschung des Vorzeichens
              \[
                  c_2(m_{12}+m_{12,3})= \frac{2y_3}{(x_3-x_1)(x_3-x_2)}\frac{(x_3-x_1)(x_3-x_2)}{2y_3}= 1
              \]
        \item Für den Fall $P_2=P_3$ und $P_1+P_2\neq P_3$ haben wir $(x_2,y_2)=(x_3,y_3)$ und dadurch $y_3=y_1+m_{12}(x_3-x_1)$ und
              \[
                  m_{12}+m_{12,3}=\frac{y_3+y_1+m_{12}(x_3-x_1)}{x_3-\hat{x}}= \frac{2y_3}{x_3-\hat{x}}
              \]
              Wir nehmen wieder die Ableitung von der Gleichung von oben und setzen jetzt $x=x_2=x_3$
              \[
                  f'(x_3)+2m_{12}y_3 = (x_3-x_1)(x_3-\hat{x})
              \]
              Dass können wir nun umformen wieder zu
              \[
                  m(P_2,P_3)-m(P_1,P_2)= \frac{(x_3-x_1)(x_3-\hat{x})}{2y_3} \Leftrightarrow
                  \frac{m(P_2,P_3)-m(P_1,P_2)}{(x_3-\hat{x})}= \frac{x_3-x_1}{2y_3} \Leftrightarrow
                  c_2 = \frac{x_3-x_1}{2y_3}
              \]
              und damit gilt dann für $y_2=y_3\neq 0$ \(c_2(m_{12}+m_{12,3})=1\)
        \item Der Fall $P_2=P_3$ und $P_1+P_2=P_3$ kann nach Lemma 2.1 nicht auftretten
    \end{enumerate}
\end{proof}

Wir haben nun die Hilfsmittel, um Lemma~\ref{lem:dreipunktsumme} zu beweisen.
\begin{proof}
    Seien $P_1,P_2,P_3\in E$ und
    \[
        P_4=(P_1+P_2)+P_3=(x_4,y_4).
    \]
    Setze $P_1+P_2=(\hat{x},\hat{y})$. Nach der Additionsformel gilt
    $\hat{x}=-x_1-x_2+m_{12}^2$ und damit
    \[
        x_4=-\hat{x}-x_3+m_{12,3}^2
        =x_1+x_2-x_3+m_{12,3}^2-m_{12}^2.
    \]
    Nach Hilfssatz~\ref{hilfs:steigungsidentitaet} gilt
    $m_{12,3}=\frac{1}{c_2}-m_{12}$, also
    \[
        \begin{aligned}
            x_4
             & =x_1+x_2-x_3+\left(\frac{1}{c_2}-m_{12}\right)^2-m_{12}^2 \\
             & =x_1+x_2-x_3+\frac{1}{c_2^2}-2\frac{m_{12}}{c_2}.
        \end{aligned}
    \]
    Vergleicht man diesen Ausdruck mit der in
    Lemma~\ref{lem:dreipunktsumme} behaupteten Formel für $x_4$,
    so ist es ausreichend, als $m_{12}= c_1+x_1c_2$ darzustellen.
    Dies folgt direkt aus den Definitionen von $c_1$ und $c_2$, denn
    \[
        \begin{aligned}
            c_1+x_1c_2
             & =\frac{x_1m_{23}-x_3m_{12}}{x_1-x_3}
            +x_1\frac{m_{12}-m_{23}}{x_1-x_3}       \\
             & =\frac{x_1m_{12}-x_3m_{12}}{x_1-x_3}
            =m_{12}.
        \end{aligned}
    \]
    Damit folgt
    \[
        \begin{aligned}
            x_4
             & =x_1+x_2-x_3+\frac{1}{c_2^2}
            -2\frac{c_1+x_1c_2}{c_2}         \\
             & =-x_1+x_2-x_3+\frac{1}{c_2^2}
            -2\frac{c_1}{c_2}.
        \end{aligned}
    \]
    Um den Ausdruck für $y_4$ zu zeigen, wird zunächst die Additionsformel
    \[
        y_4=-y_3-m_{12,3}(x_4-x_3)
    \]
    verwendet. Da $x_2\neq x_3$ gilt,
    \[
        m_{23}=\frac{y_3-y_2}{x_3-x_2},
    \]
    also
    \[
        y_3=y_2+m_{23}(x_3-x_2).
    \]
    Weiter gilt
    \begin{align*}
        m_{23}
         & =\frac{m_{23}(x_1-x_3)}{x_1-x_3}                        \\
         & =\frac{x_1m_{23}-x_3m_{12}+x_3(m_{12}-m_{23})}{x_1-x_3} \\
         & =\frac{x_1m_{23}-x_3m_{12}}{x_1-x_3}
        +x_3\frac{m_{12}-m_{23}}{x_1-x_3}                          \\
         & =c_1+x_3c_2.
    \end{align*}
    Damit folgt
    \[
        y_3=y_2+(x_3-x_2)(c_1+x_3c_2).
    \]

    Die zu zeigende Gleichung
    \[
        y_4=-y_2-(x_4-x_2)(c_1+x_4c_2)
    \]
    ist äquivalent zu
    \[
        y_4+y_2+(x_4-x_2)(c_1+x_4c_2)=0.
    \]
    Nun werden die oben erhaltenen Gleichungen eingesetzt
     \begin{align*}
         y_4=-y_3-m_{12,3}(x_4-x_3) =& -(y_2+(x_3-x_2)(c_1+x_3c_2))-m_{12,3}(x_4-x_3)\\
         =& -y_2-(x_3-x_2)(c_1+x_3c_2)-m_{12,3}(x_4-x_3)\\
     \end{align*}

     Es muss also nun gelten 
     \begin{align*}
        -(x_3-x_2)(c_1+x_3c_2)-m_{12,3}(x_4-x_3) =& (x_4-x_2)(c_1+x_4c_2)\\
        =& (x_4-x_2)(c_1+x_4c_2) + (x_3-x_2)(c_1+x_3c_2)
     \end{align*}

    Es ist $x_4 = x_1+x_2-x_3+m_{12,3}^2-m_{12}^2$

     \textbf{Ab hier quatsch}

    Weiter gilt nach der Additionsformel
    \[
        y_4=-y_3-m_{12,3}(x_4-x_3).
    \]
    Um die behauptete Darstellung von $y_4$ zu zeigen, betrachten wir daher
    \[
        y_4+y_2+(x_4-x_2)(c_1+x_4c_2).
    \]

    \[
        y_4=-y_3-m_{12,3}(x_4-x_3)
    \]



    \[
    y_3=y_2+(x_3-x_2)(c_1+x_3c_2)
\]

    Daraus können wir dann mit der Additonsformel schreiben,
    \[
        y_4=-\hat{y}-m_{12,3}(x_4-\hat{x}) = -y_2-m_{23}(\hat{x}-x_2)-m_{12,3}(x_4-\hat{x})
        = -y_2-m_{23}(\hat{x}-x_2)-m_{12,3}(-x_1+x_2-x_3+\frac{1}{c_2^2}
        -2\frac{c_1}{c_2}-\hat{x})
    \]
    als Notiz:
    \[
        \hat{y}= -y_2-m_{23}(\hat{x}-x_2)
    \]
    mit $m_{12,3}= \frac{1}{c_2}-m_{12}$
    \[
        -y_2-m_{23}(\hat{x}-x_2)-(\frac{1}{c_2}-m_{12})(x_4-\hat{x})
    \]
    \[
        -y_2-m_{23}(-x_1-x_2+m_{12}^2-x_2)-(\frac{1}{c_2}-m_{12})(x_4+x_1+x_2-m_{12}^2)
    \]
    weitere Notiz:
    \[
        \hat{x}=-x_1-x_2+m_{12}^2
    \]



\end{proof}

\begin{lemma}
    \label{lem:ass-hilfseigenschaften}
    Für alle $P,P_1,P_2\in E$ gilt,
    \begin{enumerate}[label=(\alph*)]
        \item \(-(-P)=P\)
        \item $-(P_1+P_2)=(-P_1)+(-P_2)$
        \item $(P_1+P_2)+(-P_1)=P_1+(P_2+(-P_1))=P_2$
        \item $P_1+P_2=\mathcal{O} \Rightarrow P_2=-P_1$
    \end{enumerate}
\end{lemma}

\begin{proof}
    \begin{enumerate}[label=(\alph*)]
        \item Sei $P=(x,y)\in E\setminus\{\mathcal{O}\}$, dann ist $-P=(x,-y)\in E\setminus\{\mathcal{O}\}$, also gilt $-(-P)=-(x,-y)=(x,y)=P$. Es gilt nach Definition der Identität $-\mathcal{O}=\mathcal{O}$, damit dann auch \(-(-\mathcal{O})=\mathcal{O}\)
        \item Wenn $P_1 =\mathcal{O}\lor P_2 =\mathcal{O}$ folgt direkt,
              \begin{align*}
                  P_1 =\mathcal{O}: \quad -(P_1+P_2) & =-(\mathcal{O}+P_2)=-P_2=\mathcal{O}+(-P_2) \\
                                                     & = (-\mathcal{O})+(-P_2)=(-P_1)+(-P_2)
              \end{align*}
              \begin{align*}
                  P_2 =\mathcal{O}: \quad -(P_1+P_2) & =-(P_1+\mathcal{O})=-P_1= (-P_1)+\mathcal{O} \\
                                                     & = (-P_1)+(-\mathcal{O})=(-P_1)+(-P_2)
              \end{align*}
              Nun sollen $P_1,P_2\in E \setminus\{\mathcal{O}\}$
              Wir unterscheiden 2 Fälle:
              Erster Fall $P_2=-P_1$, dann gilt $-P_2=-(-P_1)=P_1$ und damit auch
              \[
                  -(P_1+P_2)=-(P_1+(-P_1))=-\mathcal{O}=-\mathcal{O}= (-P_1)+P_1=(-P_1)+(-P_2).
              \]
              Für den anderen Fall gilt \(P_2\neq -P_1\), dann gilt per Definition der Wohldefiniertheit~\ref{lem:wohldefiniertheit}
              \[
                  P_1+P_2=(-x_1-x_2+m_{12}^2,-y_1-m_{12}(x_3-x_1)).
              \]
              Weiter gilt dann auch \(-P_2\neq-(-P_1)\) und \(-P_1=(x_1,-y_1),-P_2=(x_2,-y_2)\), dadurch gilt dann \(m(-P_1,-P_2)=-m(P_1,P_2)\)
              \begin{align*}
                  (-P_1)+(-P_2)= & (-x_1-x_2+m(P_1,P_2)^2,y_1-m(P_1,P_2)(x_3-x_1))
                  =              & -(P_1+P_2).
              \end{align*}
        \item Für $P_1=\mathcal{O}$ oder $P_2=\mathcal{O}$ folgt die Behauptung unmittelbar. Sei also $P_1,P_2\in E\setminus\{\mathcal{O}\}$. Für $P_2=-P_1$ gilt $(P_1+P_2)+(-P_1)=\mathcal{O}+P_2=P_2$. Der Fall $P_1+P_2=P_1$ kann nach Hilfssatz~\ref{hilfs:fixpunkt-addition} nicht auftreten. Sei daher $P_2\neq -P_1$ und $P_1+P_2\neq P_1$ und setze $P_3=-P_1$. Dann gilt $x_3=x_1$ und $y_3=-y_1$.

              Falls $P_1+P_2\neq P_3$, folgt wie in Hilfssatz~\ref{hilfs:steigungsidentitaet}
              \[
                  m_{12}+m_{12,3}
                  =
                  \frac{y_3+y_1+m_{12}(x_3-x_1)}{x_3-\hat{x}}
                  =0.
              \]
              Falls $P_1+P_2=P_3$, gilt entsprechend
              \[
                  m_{12}+m_{12,3}
                  =
                  \frac{(x_3-x_1)(x_3-x_2)}{2y_3}
                  =0,
              \]
              da $x_3=x_1$. Somit gilt stets $m_{12,3}=-m_{12}$.

              Mit $P_1+P_2=(\hat{x},\hat{y})$ und $\hat{x}=m_{12}^2-x_1-x_2$ folgt
              \[
                  x_4=m_{12,3}^2-\hat{x}-x_3
                  =m_{12}^2-(m_{12}^2-x_1-x_2)-x_1=x_2,
              \]
              sowie
              \[
                  y_4=-y_3-m_{12,3}(x_4-x_3)
                  =y_1+m_{12}(x_2-x_1)=y_2.
              \]
              Damit ist $P_4=P_2$ und somit $(P_1+P_2)+(-P_1)=P_2$. Ersetzen wir $P_1$ durch $-P_1$, folgt mit Teil (a) und der Kommutativität
              \[
                  P_2=((-P_1)+P_2)+(-(-P_1))
                  =P_1+(P_2+(-P_1)).
              \]
              Damit gilt
              \[
                  (P_1+P_2)+(-P_1)=P_1+(P_2+(-P_1))=P_2.
              \]
        \item Wenn $P_1+P_2=\mathcal{O}$, dann nutzen wir (c)
              \[
                  P_2=(P_1+P_2)+(-P_1)=\mathcal{O}+(-P_1)=-P_1.
              \]
    \end{enumerate}
\end{proof}

Nun können wir den Beweis der Assoziativität durchführen.
\begin{proof}[Beweis der Assoziativität]
    Wenn $P_1,P_2$ oder $P_3$ gleich $\mathcal{O}$ gilt, dann folgt direkt
    \[
        P_1=\mathcal{O}: \quad (P_1+P_2)+P_3=(\mathcal{O}+P_2)+P_3=P_2+P_3 = \mathcal{O}+(P_2+P_3)=P_1+(P_2+P_3)
    \]
    \[
        P_2=\mathcal{O}: \quad (P_1+P_2)+P_3=(P_1+\mathcal{O})+P_3 =P_1+P_3 = P_1+(\mathcal{O}+P_3)=P_1+(P_2+P_3)
    \]
    \[
        P_3=\mathcal{O}: \quad (P_1+P_2)+P_3=(P_1+P_2)+\mathcal{O} =P_1+P_2 = P_1+(P_2+\mathcal{O})=P_1+(P_2+P_3)
    \]
    Wenn nun $P_3=P_1$, dann gilt aufgrund der Kommutativität der Addition
    \[
        (P_1+P_2)+P_3=(P_1+P_2)+P_1=P_1+P_1+(P_1+P_2)=P_1+(P_2+P_3).
    \]
    Wenn nun $P_3=-P_1$, dann gilt mit der Theorem \ref{lem:ass-hilfseigenschaften} (c)
    \[
        (P_1+P_2)+P_3 = (P_1+P_2)+(-P_1)=P_1+(P_2+(-P_1))=P_1+(P_2+P_3)
    \]
    Für den Fall, dass $P_2=-P_1$ oder $P_2=-P_3$ könenn wir auch \ref{lem:ass-hilfseigenschaften} (c) nutzen
    \begin{align*}
        P_2=-P_3:\quad
        P_2=-P_1:\quad & (P_1+P_2)+P_3
        = (P_1+(-P_1))+P_3=\mathcal{O}+P_3=P_3                                       \\
                       & = P_1+(P_3+(-P_1))=P_1+((-P_1)+P_3)=P_1+(P_2+P_3)           \\
        P_2=-P_3:\quad & (P_1+P_2)+P_3
        = (P_1+(-P_3))+P_3=P_3+(P_1+(-P_3))                                          \\
                       & = P_1 = P_1+\mathcal{O}=P_3=P_1+((-P_3)+P_3)=P_1+(P_2+P_3).
    \end{align*}
    Als nächsten Fall wird $P_1+P_2=-P_3$ betrachtet, durch Anwenden von den Hilfseigenschaften~\ref{lem:ass-hilfseigenschaften} kann man herausfinden, dass
    \[
        P_2+P_3\overset{(a)}{=}P_2+(-(-P_3))=P_2+(-(P_1+P_2))\overset{(b)}{=}P_2+((-P_1)+(-P_2))=-P_1
    \]
    und kann gefolgert werden
    \[
        (P_1+P_2)+P_3=(-P_3)+P_3=\mathcal{O}=P_1+(-P_1)=P_1+(P_2+P_3).
    \]
    Für den letzten Fall wird angenommen, dass $P_1,P_2,P_3\in E\setminus{\mathcal{O}}$ mit $P_3\neq\pm P_1$, $P_2\neq -P_1$, $P_2\neq-P_3$ und $P_1+P_2\neq-P_3$. Das stimmt genau mit den Bedingungen für Lemma~\ref{lem:dreipunktsumme} überein. Zuvor wurde gezeigt, dass
    \[
        P_1+P_2=-P_3
        \quad\Longleftrightarrow\quad
        P_2+P_3=-P_1.
    \]
    Somit folgt aus $P_1+P_2\neq-P_3$ auch
    $P_2+P_3\neq-P_1$, sodass die Voraussetzungen von
    Lemma~\ref{lem:dreipunktsumme} auch nach Vertauschen von
    $P_1$ und $P_3$ erfüllt sind.  Da die dort erhaltenen Koordinaten
    von $P_4$ unter dieser Vertauschung unverändert bleiben, gilt
    \[
        (P_1+P_2)+P_3=(P_3+P_2)+P_1.
    \]
    Mit der Kommutativität der Punktaddition folgt schließlich
    \[
        (P_1+P_2)+P_3
        =(P_3+P_2)+P_1
        =P_1+(P_2+P_3).
    \]
\end{proof}

\begin{bemerkung}[Geometrische Interpretation]

\end{bemerkung}

\section{Elliptische Kurven in Charakteristik 2} (nach Rosen, S. 47ff.) \textbf{Muss man hier auch characteristik 3 ergänzen?}

Computer arbeiten tendenziell besser über Körpern der Charakteristik 2 im Vergleich zu Körpern mit großen prim Charakteristiken.Daher werden häufig Elliptische Kurven über einem Feld $\mathbb{F}_q$ mit $q=2^k$ Elementen verwendet.
(vgl. \cite{silverman06}, S.47)

Aus dieser Notwendigkeit betrachten wir nun genauer Elliptische Kurven in Charakteristik 2, da hier einige unserer vorher definierten Formeln nicht funktionieren.\\
\begin{satz}
    In Charakteristik \(2\) ist die kurze Weierstrass Gleichung \(y^2=x^3+Ax+B\) stets singulär. Daher verwendet man in Charakteristik 2 die allgemeine Weierstraß-Gleichung
\end{satz}

\[
    y^2+a_1xy+a_3y=x^3+a_2x^2+a_4x+a_6.
\]
\begin{proof}
    Betrachte
    \[
        f(x,y) = y^2-x^3-Ax-B.
    \]
    Ein Punkt auf der Kurve ist genau dann singulär, wenn beide partiellen Ableitungen verschwinden. Es gilt
    \[
        f_y=2y\equiv0,
    \] da der Körper Charakteristik 2 hat. Außerdem ist
    \[
        f_x=-3x^2-A
    \]
    Sei nun $x_0$ eine Nullstelle von $f_x$. Wir wählen ein \(y_0\) mit
    \[
        y_0^2=x_0^3+Ax_0+B
    \]
    Dann liegt $(x_0,y_0)$ auf der Kurve und es ist
    \[
        f_x(x_0,y_0)=f_y(x_0,y_0) = 0
    \]
    Damit folgt die Kurve ist singulär.
\end{proof}

Nun wollen auch für Charakterstik 2 die Gruppengesetze aufstellen.
Hierfür beginnen wir mit der allgmeinen Weierstraß-Gleichung für eine elliptische Kurv $E$
\[
    y^2+a_1xy+a_3y=x^3+a_2x^2+a_4x+a_6.
\]
Man schreibt diese in die Form
\[
    f(x,y)=y^2+a_1xy+a_3y-x^3-a_2x^2-a_4x-a_6.
\]
Dann sind die partiellen Ableitungen gegeben durch
\[
    f_x=a_1y-3x^2-2a_2x-a_4,\qquad f_y=2y+a_1x+a_3.
\]
In Charakteristik \(2\) vereinfacht sich dies zu
\[
    f_x=a_1y+x^2+a_4,\qquad f_y=a_1x+a_3.
\]
Ein Punkt \((x,y)\) ist wie bereits früher eingeführt, genau dann singulär, wenn
\[
    f(x,y)=0,\qquad f_x(x,y)=0,\qquad f_y(x,y)=0.
\]
Fall 1: $a_1\neq 0$
Man kann aus \(a_1x+a_3=0\) folgern
\[
    x=\frac{a_3}{a_1}.
\]
Eingesetzt in \(a_1y+x^2+a_4=0\) liefert es,
\[
    a_1y+\left(\frac{a_3}{a_1}\right)^2+a_4=0
\]
\[
    y= \frac{a_3^2}{a_1^3}+ \frac{a_4}{a_1}.
\]
wenn der Punkt zusätzlich \(f(x,y)=0\) erfüllt handelt es sich um einen singulären Punkt.

Fall 2: $a_1=0$
Es gilt nun,
\[
    f_x=x^2+a_4,\qquad f_y=a_3.
\]
also muss \(a_3= 0\) und $f(x,y)=0$ gelten, sodass $(x,y)$ ein singulärer Punkt ist.\\

Nun können wir die Addition von Punkten definieren. Es gilt wieder für alle Punkte $P\in E$
\[
    P+\mathcal{O}=P.
\]
\textbf{washington S.48-49 hier weiter machen}

